% !TeX root = ../main.tex
\chapter{XML 压缩存储与查询的理论基础}

\section{XML 数据模型及存储特征}

XML 文档可抽象为有序树。元素节点具有名称和子节点序列，文本节点承载字符内容；节点顺序属于文档语义的一部分。数据库中的 XML 查询需要同时处理结构信息与内容信息：结构信息包括从根到节点的路径、节点层次及其先后关系，内容信息包括文本值及其所在节点。本文方法将叶元素的完整路径作为查询和编码的基本结构单位，并保留完整节点结构流以支持结构恢复。

数据库 XML 列以逻辑 XML 值为存取对象，但物理存储还可能涉及页内表示、超大属性存储技术（The Oversized-Attribute Storage Technique，TOAST）及数据库版本相关的压缩行为[20]。因此，分析存储成本时应区分列值、表与 TOAST relation、索引以及数据库外共享字典，不能将列值大小直接等同于系统总空间。

\section{XML 压缩方法}

XML 压缩方法利用标签、路径和文本内容中的重复性降低表示开销。通用字典或统计压缩把 XML 作为字节序列处理；结构感知方法则在编码前解析标记、路径或节点事件，并据此分离或组织不同类别的数据。XMill 将 XML 数据按结构特征分组后进行压缩，体现了结构参与数据组织的思路[5]。XSC 等 XML 数据流压缩研究则从结构事件和流式编码出发，将 XML 结构处理与数据编码联系起来[4]。

这类方法的压缩对象可以包括重复标签、完整路径、文本值及树结构。若把重复路径映射到字典标识，再将标识编码为整数，可减少结构字符串的重复存储；若把相同文本值映射为 token 标识，则可在重复值较多时减少内容占用。不同方法的字典作用域、编码单位和解码粒度不同，压缩率也会受到数据重复度及结构相似度影响。结构感知压缩仍须保存足以恢复节点顺序和边界的信息，否则紧凑表示可能无法重建原始树形语义。

\section{可查询 XML 压缩方法}

可查询 XML 压缩在压缩表示的设计中考虑查询所需的信息，使查询处理有机会避免先完整恢复原文。XGrind 通过对 XML 结构和内容进行编码，研究压缩数据上的精确匹配、前缀匹配和相关查询过程[6]。XPRESS 研究 XML 数据的可查询压缩表示[7]。Buneman、Grohe 与 Koch 讨论在压缩 XML 骨架上执行路径查询的方法[8]。这些工作表明，压缩表示可以保留或构造查询可使用的结构与内容信息。

后续研究进一步把查询负载和访问粒度纳入压缩数据组织。XQueC 根据查询需求选择结构与数据容器的组织方式[9]；XSAQCT 关注 XML 压缩数据的随机访问及查询能力[10]。XQzip 利用结构索引辅助压缩 XML 查询[11]，XBzipIndex 研究压缩 XML 的导航和搜索[12]，XXS 与 SXSI 则探索压缩文档或自索引结构上的 XPath 查询[13][14]。这些方案在支持的查询语言、编码格式、索引形式及查询时解码范围方面各有特点，不能视为同一种算法或直接等价的性能基线。

从方法角度看，可查询压缩的关键不只是“压缩后能否查询”，还包括查询直接消费什么信息、候选如何定位、查询需要解码的范围，以及辅助信息由何种阶段构建。不同研究对上述问题采用了不同取舍，这为后续分析编码信息的复用方式提供了基础。

\section{XML 路径查询与索引方法}

XML 路径查询依据元素层次选择节点或判断节点存在。路径索引研究尝试将路径结构转化为可检索的键或摘要。DataGuides 使用数据结构摘要描述半结构化数据中的实际路径[15]；IndexFabric 将路径组织为索引键以支持半结构化数据检索[16]；XISS/R 研究在关系数据库环境中存储与索引 XML[17]。针对压缩 XML，XQzip、XBzipIndex、XXS 和 SXSI 等工作把结构访问与压缩表示或索引组织结合起来[11]--[14]。

路径索引能够降低查询中与目标路径无关的访问，但索引键的生成、存储和更新也构成额外成本。路径存在判断、路径对应值的提取以及跨文档内容检索并非同一操作：路径存在可由路径摘要回答；取值还需要关联文档顺序及内容记录；跨文档内容查询则需要在记录集合层面生成候选。因此，索引结构须与查询语义和数据表示相匹配。

\section{XML 内容查询与候选过滤}

内容查询可以采用精确字符串判断、词项检索或模式匹配等不同语义。若方法使用字典将文本值映射为标识，精确 token 查询可以先将查询值映射到同一标识，再在压缩内容记录或索引中定位候选。对长文本使用哈希值时，哈希相等只能说明候选可能相同，必须通过原文比较排除碰撞。

Bloom Filter 是基于位数组和多个哈希函数的集合近似结构。任一探测位为零时可以确定元素不存在；所有探测位均为一时只能判为可能存在，仍需要精确检查。它适合作为候选否定过滤器，而不能单独作为精确查询结果[22]。压缩文档中的 Bloom 摘要、关系级倒排索引和精确谓词各自作用范围不同：前者可在单文档内否定候选，倒排索引可跨文档定位候选行，最终判断仍由精确语义负责。

\section{字典编码、Varint 与 Bloom Filter}

字典编码把重复字符串映射为整数 ID。若相同字符串在指定字典作用域内复用同一 ID，内容流便可用整数标识替代重复字符串，查询输入也可转换为相同标识进行比较。收益取决于重复率、字典作用域和字典自身的存储成本。跨文档共享字典可以摊销重复字符串的空间，但字典并非零成本。

变长整数编码（Varint）根据整数大小采用不同字节数。本文实现采用 LEB128（Little Endian Base 128）风格编码：每字节低 7 位承载有效位，最高位作为延续标记。较小 ID 通常使用较少字节，同时保持顺序解码。相关整数压缩研究讨论了紧凑表示与快速访问之间的权衡[21]；Varint 是本文采用的编码选择，不是本文单独提出的整数压缩算法。

Bloom Filter 使用位数组和多个哈希探测近似判断集合成员。若任一探测位为 0，则目标一定不存在；若所有探测位均为 1，只能判为“可能存在”，后续仍需精确检查。在构建与查询使用一致参数和标识的前提下，Bloom Filter 允许假阳性而不允许假阴性[22]。本文中，字典 ID 提供可复用查询键，Varint 控制这些键在内容流中的编码成本，Bloom Filter 以 token ID 为元素构造文档级摘要，最终结果仍由内容记录确认。

\section{PostgreSQL 扩展机制与 GIN 索引}

PostgreSQL 允许扩展注册自定义数据类型、输入输出函数、操作符及索引操作类。本文原型使用自定义类型承载压缩表示，并实现相应查询函数与操作符。数据库扩展机制使特定二进制格式能够参与 SQL 查询，但正确性仍取决于类型输入输出、操作符语义及索引回调之间的一致约定。

GIN（Generalized Inverted Index，通用倒排索引）面向可由多个键表示的复合值。索引构建阶段，操作类从数据项提取键；查询阶段，将查询值转换为键并从倒排结构获取候选；一致性回调判断候选是否需要复核[19]。键抽取与索引维护本身存在成本。本文原型的 GIN 键来自压缩内容流中的 token ID；长 token 的哈希仅用于候选产生，精确字符串比较用于排除碰撞。Bloom Filter 不作为 GIN 键来源，二者不能混为同一种索引结构。

\section{本章小结}

本章介绍了理解本文方法所需的 XML 数据模型、结构感知压缩、可查询 XML 压缩、路径与内容索引、字典编码、Varint、Bloom Filter 及 PostgreSQL GIN 机制，并梳理了相关研究形成的主要技术路线。下一章不再重复介绍这些方法，而是从编码信息如何流向查询与索引、候选如何得到精确验证及其代价如何衡量等角度，分析本文研究问题的来源。
